%!TEX program = lualatex
\documentclass[ngerman,headheight=70pt]{scrartcl}
\usepackage[ngerman]{babel}
\usepackage{blindtext}
\usepackage[utf8]{luainputenc}
\usepackage{scrlayer-scrpage}
\usepackage[T1]{fontenc}
\usepackage{fontspec}
\usepackage{multicol}
\usepackage{enumitem}
\usepackage[german=quotes]{csquotes}
\usepackage[bottom=2cm,footskip=8mm]{geometry}
\usepackage{pdfpages}
\setmainfont{Source Sans Pro}
\MakeOuterQuote{"}

\setlength{\parindent}{0pt}
\setlength{\parskip}{6pt}

\pagestyle{plain.scrheadings}
\thispagestyle{scrheadings}
\lohead{\vspace*{-1.5cm}\includegraphics[scale=0.4]{../images/stupa_logo_rot.pdf}}
\rohead{\\Präsidium des\\Studierendenparlaments}

\lofoot[Universität Hamburg · Präsidium des Studierendenparlaments\\\hrule Von-Melle-Park 5 · D-20146 Hamburg]{Universität Hamburg · Präsidium des Studierendenparlaments\\\hrule Von-Melle-Park 5 · D-20146 Hamburg}
\cofoot[]{}
\cohead[\thepage]{}

\renewcommand{\thesubsection}{TOP \arabic{subsection}}
\renewcommand{\thesubsubsection}{\arabic{subsubsection}.}

\begin{document}
    UHH · StuPa-Präsidium · Von-Melle-Park 5 · D-20146 Hamburg

    %%% PAGE ONE %%%%
    \section*{Protokoll der 1. Sitzung des SWOGA vom 27. April 2016}

    \textbf{Protokoll: Jim Martens}\\
    \textbf{Ort: VMP 5 3017}\\
    \textbf{Beginn: 17.00 Uhr}\\
    \textbf{Ende: 17.34 Uhr}

    \vspace{0.5cm}
    \begin{tabular}{ll}
        Anwesend: & Melf Johannsen, Merzo Khatamov, Ramon Weilinger \\
                  & Kolja Kolb, Kevin Högy, Vincent Orth\\
                  & Artur Brückmann, Matthias Krupse (17:08) \\
                  & \\
        Entschuldigt: & Laura Franzen \\
                      &\\
        Unentschuldigt abwesend: &  \\
    \end{tabular}

    %%% Tagesordnung %%%
    \underline{Tagesordnung}
    \begin{enumerate}[label={\textbf{Top \theenumi}},leftmargin=*]
        \item Begrüßung durch ein Mitglied des Präsidiums
        \item Arbeitsplanung für die Wahlperiode 2016/2017
        \item Konstituierung des Ausschusses
        \item Verschiedenes
    \end{enumerate}

    \textit{Beschlussfähigkeit mit 7 von 9 Mitgliedern festgestellt.}

    \subsection{Begrüßung durch ein Mitglied des Präsidiums}

    Ramon begrüßt die Mitglieder. Vincent hat eine Vorstellungsrunde
    vorgeschlagen. Diese wird durchgeführt.

    \subsection{Arbeitsplanung für die Wahlperiode 2016/2017}

    Ramon erinnert daran, dass die Zusammensetzung des Ältestenrats
    Thema des Ausschusses sein sollte.

    Ramon fragt Artur, ob er noch Unterlagen von der Klausursitzung zu
    diesem Thema hat. Artur sagt, dass er nur ein Bild gemacht hat.
    Alle wollen noch einmal schauen, ob sie noch Unterlagen dazu finden
    können.

    Artur sagt, dass ein Kernkonflikt sein kann, ob im Ältestenrat
    die Mehrheitsverhältnisse des StuPa abgebildet werden sollen.
    Ramon sagt, dass RCDS grundsätzlich dafür ist, dass nicht einfach
    nur die Mehrheit abgebildet werden sollte. Allerdings sieht er
    das Problem darin, dass nur das StuPa den Ältestenrat wählen kann.

    Melf meint, dass man erst einmal Rücksprache mit den Fraktionen
    halten sollte, da interne Positionen noch nicht klar sind. Vincent
    findet es nicht so sinnvoll, erst einmal in die Gruppen zu gehen
    und dann erst im Ausschuss darüber zu sprechen. Melf sagt, dass
    es durchaus sinnvoll ist sich erst einmal im Ausschuss auszutauschen
    und dann in die Gruppen zu gehen.

    \textit{Terminfindung: Nächstes Treffen Mittwoch 11. Mai 18.30 Uhr}

    Vincent fragt, ob es weitere Themen gibt, mit denen sich der Ausschuss
    beschäftigen kann. Artur meint, dass man erst einmal das Thema
    Ältestenrat fertig machen sollte, bevor man weitere Themen festzurrt.

    Melf sagt, dass sich der Ausschuss die Vertretungsregelung ansehen
    sollte. Außerdem könnte man sich auch die Wahlzeit ansehen.

    Vincent meint, dass man sich auf jeden Fall die Fachschaftsrahmenordnung
    ansehen und eine Realitätsanpassung machen sollte.

    Es gibt keinen Widerspruch sich zunächst ausschließlich mit dem
    Ältestenrat zu befassen.

    \subsection{Konstituierung des Ausschusses}

    Vincent kandidiert für den Vorsitz. Wahl per Handzeichen.
    Vincent wird einstimmig gewählt. Er nimmt die Wahl an.

    Melf kandidiert für den stellvertretenden Vorsitz. Wahl per Handzeichen.
    Melf wird einstimmig gewählt. Er nimmt die Wahl an.

    \subsection{Verschiedenes}

    Es gibt keinen Widerspruch gegen den Vorschlag den Verteiler
    zu öffnen.

\end{document}
